% GENERATED FILE. Edit canonical lessons, then run npm run generate:books.
% canonical-source-hash: fnv1a64:03daca229037287d
% canonical-lessons: JA-C444-kokyuu-suru, JA-C444-myaku, JA-C444-ketsuatsu, JA-C444-taijuu, JA-C444-ibiki

\chapter{Breathing and Pulse}
\label{ch:ja-breathing-and-pulse}
\hlchaptermodality{444}

\begin{chapteropening}
\textbf{By the end of this chapter:} \emph{I can talk about breathing, pulse and weight.}

Complete the last lesson of chapter 444: I can talk about breathing, pulse and weight.
\end{chapteropening}

\section[\texorpdfstring{\ja{こきゅうする} (kokyū suru)}{こきゅうする (kokyū suru)}]{\ja{こきゅうする} (kokyū suru) — to breathe}

\label{lesson:JA-C444-kokyuu-suru}



\begin{quote}
Before the new one: say the Japanese for lips, then the Japanese for an eyelid.
\end{quote}

\subsection*{You'll want to know: \ja{こきゅうする}}
\textbf{\ja{こきゅうする}} — \emph{kokyū suru} — \textquotedblleft{}to breathe\textquotedblright{}.

\begin{quote}
\textbf{\ja{ゆっくり} \ja{こきゅうして} \ja{ください。}} — \emph{yukkuri kokyū shite kudasai} — \textquotedblleft{}Please breathe slowly.\textquotedblright{}
\end{quote}

Polite \textbf{\ja{こきゅうします}}, \emph{kokyūshimasu}; a \ja{する} verb. Past (polite): \textbf{\ja{かんじゃが} \ja{ゆっくり} \ja{こきゅうしました。}}, \emph{kanja ga yukkuri kokyū shimashita}.

\begin{cousinweb}
\textbf{\ja{こきゅうする}} (\emph{kokyū suru}) is Sino-Japanese (kango) \emph{kokyū} with the native verb \emph{suru}, \textquotedblleft{}to do\textquotedblright{}: two Chinese characters read with their on'yomi, readings borrowed from Chinese. They are \emph{ko}, \textquotedblleft{}call\textquotedblright{} (Middle Chinese \emph{xu}) and \emph{kyū}, \textquotedblleft{}inhale\textquotedblright{} (\emph{xip}) (KANJIDIC2; Middle Chinese in Baxter's transcription).
\end{cousinweb}

\begin{grammarlens}[title={What you've built}]
One more word for describing what you lived and felt.
\end{grammarlens}

\subsection*{Your turn}
\begin{itemize}
\raggedright
  \item \emph{Say it:} \emph{kokyū suru}
  \item \emph{Say it:} \emph{kokyū suru}, once more
  \item \emph{Say it:} \emph{yukkuri kokyū shite kudasai} — Please breathe slowly.
  \item \emph{Recall:} say the Japanese for lips, then the Japanese for an eyelid, then say \emph{kokyū suru} again
\end{itemize}

\begin{tcolorbox}[breakable,colback=teal!4,colframe=teal!35!black,title={Before you move on}]
What does \ja{こきゅうする} mean? (\textquotedblleft{}To breathe\textquotedblright{}.) Say it once more.
\end{tcolorbox}
\section[\texorpdfstring{\ja{みゃく} (myaku)}{みゃく (myaku)}]{\ja{みゃく} (myaku) — a pulse}

\label{lesson:JA-C444-myaku}



\begin{quote}
Before the new one: say the Japanese for an eyelid, then the Japanese for to breathe.
\end{quote}

\subsection*{You'll want to know: \ja{みゃく}}
\textbf{\ja{みゃく}} — \emph{myaku} — \textquotedblleft{}a pulse\textquotedblright{}.

\begin{quote}
\textbf{\ja{かんごしが} \ja{みゃくを} \ja{はかりました。}} — \emph{kangoshi ga myaku o hakarimashita} — \textquotedblleft{}The nurse took my pulse.\textquotedblright{}
\end{quote}

\begin{cousinweb}
\textbf{\ja{みゃく}} (\emph{myaku}) is Sino-Japanese (kango): one Chinese character read with its on'yomi, a reading borrowed from Chinese. It is \emph{myaku}, \textquotedblleft{}pulse\textquotedblright{} (Middle Chinese \emph{meak}) (KANJIDIC2; Middle Chinese in Baxter's transcription).
\end{cousinweb}

\begin{grammarlens}[title={What you've built}]
One more word for describing what you lived and felt.
\end{grammarlens}

\subsection*{Your turn}
\begin{itemize}
\raggedright
  \item \emph{Say it:} \emph{myaku}
  \item \emph{Say it:} \emph{myaku}, once more
  \item \emph{Say it:} \emph{kangoshi ga myaku o hakarimashita} — The nurse took my pulse.
  \item \emph{Recall:} say the Japanese for an eyelid, then the Japanese for to breathe, then say \emph{myaku} again
\end{itemize}

\begin{tcolorbox}[breakable,colback=teal!4,colframe=teal!35!black,title={Before you move on}]
What does \ja{みゃく} mean? (\textquotedblleft{}A pulse\textquotedblright{}.) Say it once more.
\end{tcolorbox}
\section[\texorpdfstring{\ja{けつあつ} (ketsuatsu)}{けつあつ (ketsuatsu)}]{\ja{けつあつ} (ketsuatsu) — blood pressure}

\label{lesson:JA-C444-ketsuatsu}



\begin{quote}
Before the new one: say the Japanese for to breathe, then the Japanese for a pulse.
\end{quote}

\subsection*{You'll want to know: \ja{けつあつ}}
\textbf{\ja{けつあつ}} — \emph{ketsuatsu} — \textquotedblleft{}blood pressure\textquotedblright{}.

\begin{quote}
\textbf{\ja{そふは} \ja{けつあつが} \ja{たかいです。}} — \emph{sofu wa ketsuatsu ga takai desu} — \textquotedblleft{}My grandfather has high blood pressure.\textquotedblright{}
\end{quote}

\begin{cousinweb}
\textbf{\ja{けつあつ}} (\emph{ketsuatsu}) is Sino-Japanese (kango): two Chinese characters read with their on'yomi, readings borrowed from Chinese. They are \emph{ketsu}, \textquotedblleft{}blood\textquotedblright{} (Middle Chinese \emph{xwet}) and \emph{atsu}, \textquotedblleft{}pressure\textquotedblright{} (\emph{'aep}) (KANJIDIC2; Middle Chinese in Baxter's transcription).
\end{cousinweb}

\begin{grammarlens}[title={What you've built}]
One more word for describing what you lived and felt.
\end{grammarlens}

\subsection*{Your turn}
\begin{itemize}
\raggedright
  \item \emph{Say it:} \emph{ketsuatsu}
  \item \emph{Say it:} \emph{ketsuatsu}, once more
  \item \emph{Say it:} \emph{sofu wa ketsuatsu ga takai desu} — My grandfather has high blood pressure.
  \item \emph{Recall:} say the Japanese for to breathe, then the Japanese for a pulse, then say \emph{ketsuatsu} again
\end{itemize}

\begin{tcolorbox}[breakable,colback=teal!4,colframe=teal!35!black,title={Before you move on}]
What does \ja{けつあつ} mean? (\textquotedblleft{}Blood pressure\textquotedblright{}.) Say it once more.
\end{tcolorbox}
\section[\texorpdfstring{\ja{たいじゅう} (taijū)}{たいじゅう (taijū)}]{\ja{たいじゅう} (taijū) — one's weight}

\label{lesson:JA-C444-taijuu}



\begin{quote}
Before the new one: say the Japanese for a pulse, then the Japanese for blood pressure.
\end{quote}

\subsection*{You'll want to know: \ja{たいじゅう}}
\textbf{\ja{たいじゅう}} — \emph{taijū} — \textquotedblleft{}one's weight\textquotedblright{}.

\begin{quote}
\textbf{\ja{たいじゅうが} \ja{すこし} \ja{ふえました。}} — \emph{taijū ga sukoshi fuemashita} — \textquotedblleft{}I have put on a little weight.\textquotedblright{}
\end{quote}

\begin{cousinweb}
\textbf{\ja{たいじゅう}} (\emph{taijū}) is Sino-Japanese (kango): two Chinese characters read with their on'yomi, readings borrowed from Chinese. They are \emph{tai}, \textquotedblleft{}body\textquotedblright{} (Middle Chinese \emph{thejX}) and \emph{jū}, \textquotedblleft{}heavy\textquotedblright{} (\emph{drjowngX}) (KANJIDIC2; Middle Chinese in Baxter's transcription).
\end{cousinweb}

\begin{grammarlens}[title={What you've built}]
One more word for describing what you lived and felt.
\end{grammarlens}

\subsection*{Your turn}
\begin{itemize}
\raggedright
  \item \emph{Say it:} \emph{taijū}
  \item \emph{Say it:} \emph{taijū}, once more
  \item \emph{Say it:} \emph{taijū ga sukoshi fuemashita} — I have put on a little weight.
  \item \emph{Recall:} say the Japanese for a pulse, then the Japanese for blood pressure, then say \emph{taijū} again
\end{itemize}

\begin{tcolorbox}[breakable,colback=teal!4,colframe=teal!35!black,title={Before you move on}]
What does \ja{たいじゅう} mean? (\textquotedblleft{}One's weight\textquotedblright{}.) Say it once more.
\end{tcolorbox}
\section[\texorpdfstring{\ja{いびき} (ibiki)}{いびき (ibiki)}]{\ja{いびき} (ibiki) — snoring; a snore}

\label{lesson:JA-C444-ibiki}



\begin{quote}
Before the new one: say the Japanese for blood pressure, then the Japanese for one's weight.
\end{quote}

\subsection*{You'll want to know: \ja{いびき}}
\textbf{\ja{いびき}} — \emph{ibiki} — \textquotedblleft{}snoring; a snore\textquotedblright{}.

\begin{quote}
\textbf{\ja{ちちの} \ja{いびきは} \ja{おおきいです。}} — \emph{chichi no ibiki wa ōkii desu} — \textquotedblleft{}My father snores loudly.\textquotedblright{}
\end{quote}

\begin{cousinweb}
\textbf{\ja{いびき}} (\emph{ibiki}) is a native Japanese word (wago), inherited rather than borrowed from Chinese; its kanji spelling was attached to it later (JMdict).
\end{cousinweb}

\begin{grammarlens}[title={What you've built}]
One more word for describing what you lived and felt.
\end{grammarlens}

\subsection*{Your turn}
\begin{itemize}
\raggedright
  \item \emph{Say it:} \emph{ibiki}
  \item \emph{Say it:} \emph{ibiki}, once more
  \item \emph{Say it:} \emph{chichi no ibiki wa ōkii desu} — My father snores loudly.
  \item \emph{Recall:} say the Japanese for blood pressure, then the Japanese for one's weight, then say \emph{ibiki} again
\end{itemize}

\begin{tcolorbox}[breakable,colback=teal!4,colframe=teal!35!black,title={Before you move on}]
What does \ja{いびき} mean? (\textquotedblleft{}Snoring; a snore\textquotedblright{}.) Say it once more.
\end{tcolorbox}
